% ==== Trustworthy AI ==========================================================
\chapter{Trustworthy AI: Safety, Security, Privacy, and Assurance}
\label{ch:trustworthy-ai}

\begin{learningobjectives}
After this chapter, you can:
\begin{itemize}
    \item distinguish reliability, functional safety, cybersecurity, privacy,
          robustness, and responsible governance;
    \item turn a use case into hazards, threats, requirements, and test
          evidence;
    \item identify attacks across the data--model--system lifecycle;
    \item explain why model confidence is not a safety argument; and
    \item draft a small assurance case with measurable release gates.
\end{itemize}
\end{learningobjectives}

\begin{plainlanguage}[The umbrella analogy]
A trustworthy AI system is like a well-run railway. \emph{Reliability} asks
whether trains usually arrive. \emph{Safety} asks whether failures can injure
people. \emph{Security} asks whether an adversary can manipulate signals.
\emph{Privacy} asks whether passenger data is misused. \emph{Governance} asks
who owns each decision and which evidence permits operation. A punctual train
is not automatically a safe, secure, or privacy-preserving one.
\end{plainlanguage}

\section{From ``accurate model'' to trustworthy system}

NIST's AI Risk Management Framework organizes voluntary risk work around four
functions: \emph{govern}, \emph{map}, \emph{measure}, and \emph{manage}.
Trustworthiness is contextual: the relevant harm, affected people, operating
conditions, and risk tolerance must be stated rather than assumed
\textsuperscript{\cite{nistAIRiskManagement2023}}.

\begin{center}
\begin{tblr}{width=\linewidth,colspec={Q[l,2.5cm]X X},hlines,
             row{1}={font=\bfseries}}
Property & Engineering question & Typical evidence \\
Reliability & Does the system keep providing its specified service? & Failure
rate, uptime, recovery tests \\
Functional safety & Can faults create unacceptable physical harm? & Hazard
analysis, safety requirements, fault injection \\
Security & Can an adversary violate confidentiality, integrity, or
availability? & Threat model, penetration tests, supply-chain controls \\
Privacy & Can information about a person be exposed or misused? & Data-flow
map, access tests, privacy analysis \\
Robustness & Does behaviour remain acceptable under bounded variation? &
Shift, stress, perturbation, and metamorphic tests \\
Fairness & Are errors and benefits distributed acceptably for the context? &
Subgroup metrics, uncertainty, impact analysis \\
Accountability & Can decisions, owners, and evidence be reconstructed? &
Versioned records, approvals, logs, incident process
\end{tblr}
\captionof{table}{Trustworthiness properties connect distinct engineering
questions to distinct forms of evidence.}
\label{tab:ch08-functional-security-inline-01}
\end{center}

For hazards $h_i$, a first screening model is
\begin{equation}
    R = \sum_i P(h_i \mid c)\,C(h_i \mid c),
\end{equation}
where $c$ describes the operational context and $C$ is consequence severity.
The equation disciplines thought; it does not manufacture trustworthy numbers.
Rare catastrophic events, uncertain probabilities, distribution shift, and
unequal effects require scenarios and uncertainty ranges alongside a scalar.

\begin{engineernote}
Model accuracy is a component metric. System risk also depends on exposure,
human interaction, sensors, software, actuators, access rights, fallback
behaviour, and the severity of being wrong. A 99.9\% accurate classifier may
still be unacceptable if its remaining errors trigger hazardous actions.
\end{engineernote}

\section{The assurance loop}

\begin{figure}[ht]
\centering
\begin{tikzpicture}[
    node distance=5mm and 7mm,
    every node/.style={font=\sffamily\scriptsize,align=center},
    assurance step/.style={draw,rounded corners,fill=LimeGreen!8,
                           minimum width=2.25cm,minimum height=8mm},
    arr/.style={-{Stealth[length=2mm]},thick}]
\node[assurance step] (scope) {Scope\\system and context};
\node[assurance step,right=of scope] (analyse) {Analyse\\hazards and threats};
\node[assurance step,right=of analyse] (require) {Specify\\claims and limits};
\node[assurance step,below=of require] (control)
    {Control\\prevent, detect, recover};
\node[assurance step,left=of control] (verify) {Verify\\tests and evidence};
\node[assurance step,left=of verify] (operate) {Operate\\monitor and respond};
\draw[arr] (scope) -- (analyse);
\draw[arr] (analyse) -- (require);
\draw[arr] (require) -- (control);
\draw[arr] (control) -- (verify);
\draw[arr] (verify) -- (operate);
\draw[arr] (operate) -- (scope);
\end{tikzpicture}
\caption{Assurance is a closed lifecycle, not a certificate attached at the
end.}
\label{fig:assurance-loop}
\end{figure}

\subsection{Define the system boundary}

Begin with the complete sociotechnical system: users, model, training and
evaluation data, APIs, retrieval stores, tools, sensors, actuators, operators,
and downstream consumers. State the intended purpose and explicit exclusions.
Then describe an \keyterm{operational design domain}: conditions under which
the evidence is meant to hold.

\subsection{Hazard analysis and risk assessment}

A \keyterm{hazard} is a condition that can lead to harm; a \keyterm{threat} is
a circumstance or actor capable of exploiting a vulnerability. Hazard analysis
asks ``what can hurt whom?'' Threat modelling adds ``who or what can make it
happen, through which path?''

\begin{enumerate}
    \item Name assets and affected stakeholders.
    \item Describe losses or harms before guessing causes.
    \item Identify hazardous system states and initiating events.
    \item Estimate severity, exposure, controllability, and uncertainty.
    \item Derive testable safety and security requirements.
    \item Allocate controls and owners; document residual risk.
\end{enumerate}

In automotive engineering, ISO 26262 structures functional-safety work across
the lifecycle and uses Automotive Safety Integrity Levels (ASILs). The level is
derived by the standard's method; it must not be guessed from model accuracy
\textsuperscript{\cite{isoRoadVehiclesFunctional2018}}. ISO/SAE 21434 addresses
vehicle cybersecurity engineering across its lifecycle
\textsuperscript{\cite{isoRoadVehiclesCybersecurity2021}}. They are
complementary scopes, not interchangeable badges.

\subsection{FTA and FMEA: top-down and bottom-up views}

\keyterm{Fault Tree Analysis (FTA)} starts with an unwanted top event and
decomposes combinations of causes. If independent rare causes $A$ and $B$ feed
an OR gate, $P(A\cup B)=P(A)+P(B)-P(A)P(B)$. For an AND gate,
$P(A\cap B)=P(A)P(B)$. Independence is an assumption to justify; shared data,
power, clocks, or software often create common-cause failures.

\keyterm{Failure Modes and Effects Analysis (FMEA)} moves bottom-up: for every
component or function, ask how it can fail, what local and system effects
follow, how the failure is detected, and what prevents or contains it. For ML,
failure modes include stale features, corrupted labels, silent schema changes,
out-of-domain inputs, unstable preprocessing, and unsafe downstream action.

For conventional random hardware failures, two useful ratios are
\begin{align}
  \mathrm{DC} &= \frac{\lambda_{DD}}{\lambda_{DD}+\lambda_{DU}}, &
  \mathrm{SFF} &= \frac{\lambda_S+\lambda_{DD}}
  {\lambda_S+\lambda_{DD}+\lambda_{DU}},
\end{align}
where $\lambda_S$ is the safe failure rate and $\lambda_{DD},\lambda_{DU}$ are
detected and undetected dangerous failure rates. These hardware-oriented
quantities do not turn systematic specification, data, or software errors into
random failures. Use the right model for the failure mechanism.

\section{ML-specific failure modes}

\begin{center}
\begin{tblr}{width=\linewidth,colspec={Q[l,2.6cm]X X},hlines,
             row{1}={font=\bfseries}}
Failure family & Example & Evidence or control \\
Specification & Proxy objective rewards the wrong behaviour & Scenario review,
stakeholder sign-off, counterexamples \\
Data & Missing subgroup, label error, leakage, stale source & Datasheet,
provenance, slice tests, schema gates \\
Statistical & Shift, poor calibration, irreducible overlap & OOD tests,
calibration plot, abstention policy \\
Implementation & Training/serving skew, numerical overflow & Golden examples,
parity tests, bounds checks \\
Interaction & Automation bias or confusing handover & Human-factors study,
usable override, training \\
Operations & Drift, dependency compromise, telemetry gap & Signed artifacts,
monitoring, rollback drill
\end{tblr}
\captionof{table}{ML failures arise at specification, data, statistical,
implementation, interaction, and operational boundaries.}
\label{tab:ch08-functional-security-inline-02}
\end{center}

If a model emits confidence $s(x)$, that number is not automatically
$P(Y=\hat y\mid x)$. Calibration must be measured on relevant data, including
important slices. A selective system may abstain when confidence is low. With
accepted set $A_\tau=\{x:s(x)\geq\tau\}$,
\begin{align}
  \mathrm{coverage}(\tau) &= \frac{|A_\tau|}{n}, &
  \mathrm{risk}(\tau) &= \frac{1}{|A_\tau|}
       \sum_{x_i\in A_\tau}\ell(f(x_i),y_i).
\end{align}
The risk--coverage curve makes the trade-off visible. Abstention is useful only
when a safe fallback exists and the operator can recognize the handover.

\begin{cautionbox}[The confidence trap]
Softmax can be extremely confident far from the training distribution. Never
write ``the model knows when it is wrong'' without a defined uncertainty
method, representative evaluation, thresholds, and a tested fallback path.
\end{cautionbox}

\section{Adversarial machine learning}

NIST classifies adversarial ML by system type, lifecycle stage, attacker goal,
capability, access, and knowledge. Its 2025 taxonomy covers predictive and
generative systems and stresses that mitigations have limits
\textsuperscript{\cite{vassilevAdversarialMachineLearning2025}}.

\begin{center}
\begin{tblr}{width=\linewidth,
             colspec={Q[l,2.3cm]Q[l,2.2cm]X X},hlines,
             row{1}={font=\bfseries}}
Attack & Stage & Engineering example & Defence-in-depth evidence \\
Poisoning/backdoor & Train or adapt & Malicious examples create a hidden trigger
& Provenance, access control, anomaly review, clean holdout \\
Evasion & Inference & Crafted input crosses a decision boundary &
Threat-specific
stress tests, input constraints, detection, safe response \\
Model extraction & Serving & Queries reconstruct behaviour or parameters & Rate
limits, output minimization, monitoring, contractual controls \\
Membership/privacy & Train/serve & Output reveals whether a record was present &
Data minimization, privacy test, regularization or formal privacy \\
Prompt injection & Application & Untrusted text competes with system
instruction &
Treat content as data, isolate tools, least privilege, confirmation \\
Tool misuse & Application & Model invokes a damaging action & Allow-list,
argument validation, sandbox, authorization outside the model \\
Supply-chain attack & Build/deploy & Dataset, package, model, or adapter is
replaced
& Pinning, hashes/signatures, scanning, bill of materials, rollback
\end{tblr}
\captionof{table}{Adversarial-ML controls must address attacks across the
training, serving, application, and supply-chain lifecycle.}
\label{tab:ch08-functional-security-inline-03}
\end{center}

\begin{engineernote}[The model is not the security boundary]
A prompt such as ``never reveal secrets'' is guidance to a probabilistic
component, not access control. Credentials, permissions, validation, spending
limits, and destructive-action approval belong in deterministic systems around
the model. Assume retrieved documents and tool output may be hostile.
\end{engineernote}

For generative systems, also test confabulation, harmful or illegal assistance,
sensitive-data disclosure, bias, information integrity, human over-reliance,
and environmental or intellectual-property risks. NIST's Generative AI Profile
maps such risks to govern--map--measure--manage actions
\textsuperscript{\cite{nistGenerativeAIProfile2024}}.

\section{Privacy by engineering, not by wishful thinking}

Start with minimization: collect only justified fields, separate identifiers,
limit retention, encrypt data, enforce least privilege, and log access. Simple
de-identification is often reversible when auxiliary information exists.

Two datasets $D,D'$ are \keyterm{adjacent} if they differ in one person's
contribution. A randomized mechanism $M$ is
$(\varepsilon,\delta)$-differentially
private if for every measurable output set $S$,
\begin{equation}
  P[M(D)\in S] \leq e^{\varepsilon}P[M(D')\in S]+\delta.
\end{equation}
This bounds how much one contribution can change the released distribution.
Smaller $\varepsilon$ generally means a stronger bound but more noise; repeated
releases consume privacy budget. The guarantee depends on adjacency,
accounting, implementation, and parameter choices---not merely on adding some
noise\textsuperscript{\cite{nearDifferentialPrivacyGuarantees2025}}.

Federated learning keeps raw training records local, but updates can still leak
information and a coordinator can still be attacked. It is an architecture,
not a proof of privacy. Combine it with explicit threat analysis and, when
appropriate, secure aggregation or differential privacy.

\section{Safety cases and governance}

An \keyterm{assurance case} is a structured argument connecting claims to
evidence under stated assumptions. It is stronger than a folder full of test
reports because it explains why each report matters.

\begin{workedexample}[A small assurance argument]
\textbf{Claim:} the vision assistant cannot directly command hazardous motion.
\begin{itemize}
    \item \textbf{Argument:} commands pass through an independent, bounded
          safety controller; uncertain perception triggers a controlled stop.
    \item \textbf{Evidence:} interface specification, static analysis,
          adversarial/OOD tests, latency bounds, fault injection, and a physical
          emergency-stop test.
    \item \textbf{Assumptions:} sensor range, braking capability, speed
          envelope,
          operator response, and maintenance interval.
    \item \textbf{Defeaters:} shared power failure, stale frames, controller
          bypass, or a test distribution that omits poor weather.
\end{itemize}
The defeaters are not embarrassing footnotes; they are the shopping list for
the next tests.
\end{workedexample}

ISO/IEC 42001 specifies requirements for establishing and continually
improving an organizational AI management system. It provides governance
structure; it does not certify that every model output is safe
\textsuperscript{\cite{isoAIManagementSystems2023}}. Applicable law and sector
standards still need specialist review; see the legal chapter for orientation,
not legal advice.

\section{Practical lab: threat model and release gate}

\begin{practicallab}[Red-team a small AI-assisted control application]
\textbf{Scenario.} A language model translates a technician's request into a
proposed maintenance action. A deterministic service validates and executes
approved actions in a simulator. The model never receives production secrets.

\textbf{Tasks.}
\begin{enumerate}
    \item Draw trust boundaries and data flows. Mark users, model host,
          retrieved manuals, tool service, logs, and simulated plant.
    \item Name at least three harms, three accidental failure modes, and three
          adversarial paths. State attacker access and capability.
    \item Write controls at prevention, detection, response, and recovery
          layers. Give each an owner and observable signal.
    \item Build a private evaluation set containing valid requests, malformed
          arguments, prompt injection inside a manual, privilege escalation,
          stale documentation, ambiguous units, and out-of-scope commands.
    \item Define release gates before testing: zero unauthorized executions;
          100\% rejection of destructive unconfirmed actions; bounded p95
          latency; logged rationale and artifact versions for every attempt.
    \item Run the suite on a baseline and one improved design. Report every
          failed case; do not replace failures after seeing results.
    \item Demonstrate rollback and write a one-page assurance case with claims,
          evidence, assumptions, residual risks, and expiry date.
\end{enumerate}

\textbf{Deliverables.} Threat-model diagram, machine-readable test cases,
results with confidence intervals where sampling is involved, incident/rollback
record, and the assurance case. A release is ``not ready'' when evidence is
missing---even if the demo is charming.
\end{practicallab}

\begin{checkpoint}
\begin{enumerate}
    \item Why can high held-out accuracy coexist with unacceptable system risk?
    \item Give one common-cause failure that invalidates an independent-events
          fault-tree calculation.
    \item What security property cannot be enforced by a prompt alone?
    \item Which assumptions define the meaning of an $(\varepsilon,\delta)$
          privacy guarantee?
    \item Write one claim--evidence--assumption triple for your own project.
\end{enumerate}
\end{checkpoint}
